\documentclass[../../main.tex]{subfiles}
\graphicspath{{../../image/}{../../figure/}}
\begin{document}

\section{熵--熵流对与速矢变换}\label{section:ch12-熵熵流对与速矢变换}

如原书第 7.4 节所述，我们的系统配有一个丰富的熵--熵流对族 $(\eta,q)$，它们可通过求解原书系统 (7.4.12) 而确定为 Riemann 不变量 $(z,w)$ 的函数，即
\begin{align}
q_z=\lambda\eta_z,\qquad q_w=\mu\eta_w.
\label{align:ch12-ent流}
\end{align}
可积性条件原书 (7.4.13) 现取形式
\begin{align}
\eta_{zw}+\frac{\lambda_w}{\lambda-\mu}\eta_z+\frac{\mu_z}{\mu-\lambda}\eta_w=0.
\label{align:ch12-熵可积}
\end{align}
当原书不等式 (7.4.16) 成立时，熵 $\eta(z,w)$ 将是原状态变量 $U$ 的凸函数，即
\begin{align}
\begin{cases}
\eta_{zz}+(R^{\top}\mathrm{D}^2_z R)\eta_z+(R^{\top}\mathrm{D}^2_w R)\eta_w\geqslant 0,\\
\eta_{ww}+(S^{\top}\mathrm{D}^2_z S)\eta_z+(S^{\top}\mathrm{D}^2_w S)\eta_w\geqslant 0.
\end{cases}
\label{align:ch12-凸性}
\end{align}

在研究过程中，我们将需要在给定边值条件下通过求解\eqref{align:ch12-ent流}或\eqref{align:ch12-熵可积}来构造具有指定形态的熵--熵流对。为验证所构造的熵满足凸性条件\eqref{align:ch12-凸性}，通常需要以一阶导数 $\eta_z$、$\eta_w$ 来估计二阶导数 $\eta_{zz}$、$\eta_{ww}$。为此目的，可以使用将\eqref{align:ch12-熵可积}分别对 $z$ 和 $w$ 求导所得的方程：
\begin{align}
\begin{cases}
\eta_{zzw}+\dfrac{\lambda_w}{\lambda-\mu}\eta_{zz}
=\dfrac{(\mu-\lambda)\lambda_{zw}+\lambda_z\lambda_w-2\lambda_w\mu_z}{(\lambda-\mu)^2}\eta_z
+\dfrac{(\lambda-\mu)\mu_{zz}-\lambda_z\mu_z+2\mu_z^2}{(\lambda-\mu)^2}\eta_w,\\[2mm]
\eta_{wwz}+\dfrac{\mu_z}{\mu-\lambda}\eta_{ww}
=\dfrac{(\mu-\lambda)\mu_{ww}-\lambda_w\mu_w+2\lambda_w^2}{(\mu-\lambda)^2}\eta_z
+\dfrac{(\lambda-\mu)\mu_{zw}+\mu_z\mu_w-2\lambda_w\mu_z}{(\mu-\lambda)^2}\eta_w.
\end{cases}
\label{align:ch12-熵二阶导数}
\end{align}

作为说明，我们考虑重要的 Lax 熵--熵流对族：
\begin{align}
\begin{cases}
\eta(z,w)=e^{kz}\left[\phi(z,w)+\dfrac{1}{k}\chi(z,w)+O\!\left(\dfrac{1}{k^2}\right)\right],\\[1mm]
q(z,w)=e^{kz}\lambda(z,w)\left[\psi(z,w)+\dfrac{1}{k}\theta(z,w)+O\!\left(\dfrac{1}{k^2}\right)\right],
\end{cases}
\label{align:ch12-lax族1}
\end{align}
\begin{align}
\begin{cases}
\eta(z,w)=e^{kw}\left[\alpha(z,w)+\dfrac{1}{k}\beta(z,w)+O\!\left(\dfrac{1}{k^2}\right)\right],\\[1mm]
q(z,w)=e^{kw}\mu(z,w)\left[\gamma(z,w)+\dfrac{1}{k}\delta(z,w)+O\!\left(\dfrac{1}{k^2}\right)\right],
\end{cases}
\label{align:ch12-lax族2}
\end{align}
其中 $k$ 是参数。它们被设计为随两个 Riemann 不变量之一剧烈变化，以便用于退耦两个特征族。为构造它们，将\eqref{align:ch12-lax族1}或\eqref{align:ch12-lax族2}中的 $\eta$ 与 $q$ 代入系统\eqref{align:ch12-ent流}，由此导出系数的递推关系，然后证明余项为 $O(k^{-2})$。族\eqref{align:ch12-lax族1}系数的递推关系如下：
\begin{align}
\psi=\phi,
\label{align:ch12-递推1}
\end{align}
\begin{align}
\lambda\theta+(\lambda\psi)_z=\lambda\chi+\lambda\phi_z,
\label{align:ch12-递推2}
\end{align}
\begin{align}
(\lambda\psi)_w=\mu\phi_w.
\label{align:ch12-递推3}
\end{align}
将\eqref{align:ch12-递推1}与\eqref{align:ch12-递推3}结合得
\begin{align}
(\mu-\lambda)\phi_w=\lambda_w\phi,
\label{align:ch12-phi方程}
\end{align}
可通过选取
\begin{align}
\phi(z,w)=\exp\int_0^w\frac{\lambda_w(z,\omega)}{\mu(z,\omega)-\lambda(z,\omega)}\,\mathrm{d}\omega
\label{align:ch12-phi显式}
\end{align}
来满足。特别地，该 $\phi$ 为正，且在紧集上一致远离零。因此，当 $k$ 充分大时，不等式\eqref{align:ch12-凸性}成立，其中第二个不等式由\eqref{align:ch12-相互作用}保证。于是，当 $k$ 很大时，Lax 熵是 $U$ 的凸函数。

\eqref{align:ch12-递推1}与\eqref{align:ch12-递推2}的重要推论是估计
\begin{align}
q-\lambda\eta=\frac{1}{k}e^{kz}\left[-\lambda_z\phi+O\!\left(\frac{1}{k}\right)\right],
\label{align:ch12-估计1}
\end{align}
\begin{align}
q-(\lambda+\varepsilon)\eta=-e^{kz}\left[\varepsilon\phi+O\!\left(\frac{1}{k}\right)\right],
\label{align:ch12-估计2}
\end{align}
其用处将在后续显现。

映射 $(z,w)\mapsto(\eta,q)$ 将成对的 Riemann 不变量带入熵--熵流对，它与由下述过程构造的速矢变换 $(z,w)\mapsto(x,t)$ 之间存在奇特的形式类比：设 $z(x,t)$ 与 $w(x,t)$ 是系统\eqref{align:ch12-系统}在 $x$-$t$ 平面某区域 $\mathcal{D}$ 上的一个 $C^1$ 解的 Riemann 不变量。在 $\mathcal{D}$ 中任一使 Jacobi 行列式 $J=z_xw_t-w_xz_t$ 不消失的点附近，映射 $(x,t)\mapsto(z,w)$ 有 $C^1$ 逆映射 $(z,w)\mapsto(x,t)$，其偏导数为 $x_z=J^{-1}w_t$、$t_z=-J^{-1}w_x$、$x_w=-J^{-1}z_t$、$t_w=J^{-1}z_x$。由于在 $\mathcal{D}$ 上 $z_t+\lambda z_x=0$、$w_t+\mu w_x=0$，推得 $J=(\lambda-\mu)z_xw_x$ 以及
\begin{align}
x_z=\mu t_z,\qquad x_w=\lambda t_w,
\label{align:ch12-速矢}
\end{align}
这应与\eqref{align:ch12-ent流}对照比较。在\eqref{align:ch12-速矢}的两个方程之间消去 $x$，得
\begin{align}
t_{zw}+\frac{\mu_w}{\mu-\lambda}t_z+\frac{\lambda_z}{\lambda-\mu}t_w=0,
\label{align:ch12-速矢方程}
\end{align}
即\eqref{align:ch12-熵可积}的类比。\footnote{在原书系统 (7.1.11) 中，\eqref{align:ch12-熵可积}与\eqref{align:ch12-速矢方程}重合，因为 $\lambda=-\mu$。}
由此，人们可通过求解线性系统\eqref{align:ch12-速矢}，或等价地线性二阶双曲方程\eqref{align:ch12-速矢方程}，来构造两方程非线性系统\eqref{align:ch12-系统}的（经典）解。等熵气体动力学系统以及数学物理中出现的其他两方程守恒律组的许多重要特解，都是通过这一过程导出的。

\end{document}
